% Copyright (c) 2026 Geovana Neves. Licensed under CC BY 4.0 (https://creativecommons.org/licenses/by/4.0/). See LICENSE-AND-AUTHORSHIP.md.
%
% fsi/text/05-the-row.tex

\section{From the row to the solver}

\begin{frame}{What a coupled row stages}
\small A row that states \code{FSI: f001} on a workflow that allows it
prepares the driver's complete input interface. Planning derives the nodes
and the maps from the effective configuration and resolves the boundaries and
the frames of each blade or of the wing; the run writes and hashes:
\par\vspace{1.5mm}
{\ftstablesize\renewcommand{\arraystretch}{1.15}
\begin{tabularx}{\linewidth}{@{}P{0.27\linewidth}R@{}}
\toprule
\textbf{File, in the point's directory} & \textbf{What it is} \\
\midrule
\code{config.json}, \code{fsi-provenance.json} & the effective configuration, and where every number came from \\
\code{fsi\_nodes.csv} & the structural nodes, inside the component, in its own frame; one identical file imported once per blade \\
\code{fsi\_node\_map.json} & the node ordering map (\code{node\_map\_file}), from the same generator \\
\code{fsi\_family\_map.json} & the emitted section-load order: which rows of the loads export belong to which blade \\
\code{fsi\_post.txt} & the native loads-export callback: update the sections, compute their loads in Newtons, export them \\
\code{fsi\_callback.py} & the Python callback: one \code{pyfs-fsi} step in the point's directory \\
\bottomrule
\end{tabularx}}
\par\vspace{1.5mm}
{\scriptsize The callback runs under the interpreter that prepared the run
(its \code{pythonw.exe} sibling on Windows): the package, with its
\code{[fsi]} extra, must be installed for that interpreter on the machine
that executes the point.\par}
\fsisource{\code{docs/fsi-workspace.md}, ``Automatic workspace coupling'';
  \code{cases/fsi\_workspace.py}, \code{wire\_workspace\_fsi};
  \code{fsi/beam.py} (PyNite from the \code{[fsi]} extra)}
\end{frame}

% RE-CHECK after the FSI branch (FSI-1, FSI-GUARD, FSI-G) merges: the
% command list and its order are written from the requirements (surface by
% boundary ID, multiquadric kernel, steady and unsteady coupling), not read
% from the merged builder; compare with a planned script of each workflow.
\begin{frame}[fragile]{The native commands, in the order they are emitted}
\relax % a body opening with a brace group would be read as the frame subtitle
{\ftstablesize\renewcommand{\arraystretch}{0.97}
\begin{tabularx}{\linewidth}{@{}P{0.47\linewidth}R@{}}
\toprule
\textbf{Command} & \textbf{What the row gives it} \\
\midrule
\code{DELETE\_AEROELASTIC\_STRUCTURAL\_NODES} & nothing: a clean start \\
\code{AEROELASTIC\_RBF\_TYPE} & \code{MULTI\_QUADRATIC}, for a beam line of nodes \\
\code{ASSIGN\_AEROELASTIC\_SURFACES} & the coupled boundaries, by their identifiers in the mesh \\
\code{ASSIGN\_AEROELASTIC\_COORDINATE\_SYSTEMS} & each blade's own frame, or the wing's \\
\code{IMPORT\_AEROELASTIC\_STRUCTURAL\_NODES} & per blade, in blade order: its frame, \code{DISABLE}, \code{fsi\_nodes.csv} \\
\code{SET\_AEROELASTIC\_WORKING\_DIRECTORY} & \code{.}, the point's directory \\
\code{SET\_AEROELASTIC\_POST\_PROCESSING\_SCRIPT} & \code{fsi\_post.txt} \\
\code{SET\_AEROELASTIC\_STRUCTURAL\_EXECUTION\_COMMAND} & the interpreter and \code{fsi\_callback.py} \\
\code{SET\_AEROELASTIC\_ITERATIONS} & \code{1}, one call per coupling step \\
\code{EXECUTE\_AEROELASTIC\_ANALYSIS} & a \code{steady} or \code{qsteady\_rotor} row: the steady coupled solve \\
\code{SET\_AEROELASTIC\_COUPLING\_IN\_UNSTEADY} & an \code{unsteady} row: \code{ENABLE}, before \code{START\_SOLVER} \\
\bottomrule
\end{tabularx}}
\par\vspace{1mm}
\begin{columns}[T]
\begin{column}{0.44\textwidth}
\scriptsize \code{fsi\_callback.py}, whole:
\begin{lstlisting}[style=ftsbase, language=Python, basicstyle=\ttfamily\tiny]
from pyflightstream.fsi.cli import main
raise SystemExit(main(['step', '--dir', '.']))
\end{lstlisting}
\end{column}
\begin{column}{0.52\textwidth}
\begin{alertblock}{Documented, not verified}
\scriptsize The command database records every command of the family as
documented, none as verified: its 26.124 rows are carried from the 26.123
manual. The coupled runs of this version are measurements, not promotions.
\end{alertblock}
\end{column}
\end{columns}
\fsidecided{\code{cases/fsi\_workspace.py}, \code{wire\_workspace\_fsi};
  \code{commands/aeroelastic\_coupling.yaml}}
\end{frame}

\begin{frame}[fragile]{The pproc sections a coupled row needs}
\begin{columns}[T]
\begin{column}{0.44\textwidth}
\small A one-blade rotor whose family is \code{Blade1}:
\begin{lstlisting}[style=ftsbase, basicstyle=\ttfamily\scriptsize]
[sections]
count = 50
include_symmetry = false

[[sections.distributions]]
families = ["Blade1"]
frame = "LOCAL_AXIS"
planes = ["XY"]
\end{lstlisting}
{\scriptsize No separate \code{[settings.aeroelastic]} table is
needed.\par}
\end{column}
\begin{column}{0.52\textwidth}
\begin{itemize}\small
  \item Exactly one XY distribution per blade, emitted in blade order, each in
        that blade's own frame (span along Z); a wing takes one along its
        span.
  \item Several blades: select the rotor's blades so \code{LOCAL\_AXIS}
        expands them in the declared blade order. The structural
        \code{blade\_count} counts every explicit blade.
  \item Duplicate cuts, aggregated families, other planes or additional
        distributions are refused, never attributed to the first matching
        blade.
\end{itemize}
\begin{takeaway}
The loads export is one flat table in creation order; the family map is what
tells the driver which rows are which blade.
\end{takeaway}
\end{column}
\end{columns}
\fsisource{\code{docs/fsi-workspace.md}, ``Automatic workspace coupling'';
  FSI tutorial, ``loads''}
\end{frame}
